\begin{table*}[t]
  \centering
  \small
  \begin{tabular}{@{}L{0.15\textwidth}L{0.22\textwidth}
    L{0.25\textwidth}L{0.29\textwidth}@{}}
    \toprule
    Exercise & Evaluated unit and support & Baseline and reference &
    Sensitivity \\
    \midrule
    Volatility
    & 49 common months, 675,234 eligible articles, four publishers
    & VIX $R^2$ 0.515 vs. 0.647, gap $-0.132$ [$-0.215$, $+0.110$]
    & Gap $-0.041$ after dropping March--April 2020 from both fits \\
    Bank network
    & Distinct unordered bank pairs in each earnings week, 33 of 34
    reference nodes mapped
    & 177 to 242 pairs vs. 210 to 172 in the July 2023 manuscript
    & Reuters only, 50 to 35. Matched mean increases $+56.3$ [5, 83]
    on six weeks and
    $+19.1$ [$-25$, 59] on fourteen weeks \\
    Overnight drift
    & 277 evaluated sessions, 42 with both legs and 123 with neither
    & $-10.15$ bp [$-25.2$, $+4.2$] vs. $+11$ bp for the reference
    MiniLM headline sample
    & Recovered clocks give 340 evaluated sessions, 110 with both
    legs, and $-7.73$ bp after refitting \\
    \bottomrule
  \end{tabular}
  \caption{The financial applications use monthly regressions,
    earnings-week networks and open-to-close portfolio returns,
    respectively. Volatility is contemporaneous fit with a paired
    six-month block interval. Bank ranges are the central 95\% of
    matched-story draws, with 682 stories per week in the six-week
    analysis and 158 in the fourteen-week analysis. Drift is gross basis
    points per evaluated session, with a 21-session block interval
  conditional on the saved predictions.}
  \label{tab:studies}
\end{table*}

\section{Financial Applications}
\label{sec:results}

Table~\ref{tab:studies} compares the measurements supported by the
three study samples and shows how they change under alternative
publisher and timing rules.

\subsection{Publisher-Month Volatility}

We apply the monthly specifications of \citet{baker2025ms,baker2026}
after reproducing the six published columns to their displayed precision
on the authors' archived data \citep{baker2025archive}. Our tracker
instead uses 675,234 consistently dated canonical articles from
Benzinga, Forbes, MarketWatch and the \emph{Wall Street Journal}.
Only the last appears in the authors' eleven-newspaper basket.
We count an article as a hit when its title and body together contain
terms from all three dictionary groups for the economy, equity markets
and volatility. We divide the hit share for each publisher by its
dispersion over 40 calibration months through December 2022, average
the standardized shares, and rescale the average to the
calibration-period mean VIX. For September 2019 through September
2023, each tracker enters its own regression of monthly VIX with an
intercept and the same outcomes from FRED. The captured-article tracker
has an $R^2$ of 0.515 against 0.647 for the authors' tracker, a difference
of $-0.132$ (paired 95\% interval $[-0.215, 0.110]$).

Omitting March and April 2020 from both regressions narrows the VIX
difference to $-0.041$. The choice of publishers also affects the
comparison with calibration held fixed. For example, the shortfall
is larger when Forbes is omitted, both with and without the two
months. In the full calendar, that omission changes the difference from
$-0.132$ to $-0.360$. Both trackers also explain less variation when
March and April 2020 are omitted, which shows how much the fit depends
on those two months. These comparisons describe historical fit on a
common calendar. Because the publisher basket determines which stories
contribute to the index, a change in source can alter more than the
volume of articles a researcher sees. It can also alter which firms
appear together in the retained stories.

\subsection{Bank Pairs in Earnings Weeks}

We compare co-mentions in the weeks of 13--19 January and 13--19 April
2020 against the July 2023 benchmark of
\citet{kazinnik2023ms,kazinnik2026}. The archived corpus supplies
links for 33 of the reference design's 34 bank nodes, in some cases
through a foreign group's listed parent, and substitutes these links
for the Factiva tags. The reconstructed networks recover a complete core
among the six largest institutions. The number of connected pairs rises
from 177 to 242, whereas the July 2023 manuscript reports a decline
from 210 to 172. Restricting CCFinance to Reuters hosts yields 50 to 35.
The comparison therefore depends on which bank relationships a
publisher mix brings into the week.

To examine that dependence, we collapse syndicated copies and compare
weeks after equating the number of stories from each publisher. Over
six weeks with 682 matched stories per week, the mean January--April
2020 increase is 56.3 connected pairs, with a central 95\% subsampling
range of [5, 83]. Extending the calendar to fourteen weeks restricts the
roster to 32 mapped bank nodes and requires common publisher quotas on
all fourteen weeks. The narrower common support retains 158 stories
per week and excludes Nasdaq, which supplied a substantial share of
the earlier matched sample.
The increase becomes 19.1 pairs [$-25$, 59]. The ranges describe
variation across draws of matched stories. The seven nonoverlapping
week comparisons in the extended sample give limited historical
evidence about how exceptional the increase was.

\subsection{Issuer-Session Drift}

The overnight analogue of \citet{lopez2025v6,lopez2026} begins with
996 captured bodies from Business Standard and Mint Lounge linked
to 171 issuers. It uses their EmbeddingGemma vectors and Yahoo price
histories in place of the authors' RavenPack-matched headlines and
CRSP prices. Following the authors' archive, we fit separate rolling
logistic regressions using prior return signs from six-month windows
\citep{lopez2026archive}. The earliest eligible predicted article for a
security supplies its session signal, and each qualifying portfolio leg
weights the securities within that leg equally. The session spread
subtracts the shorted securities' mean return from the long return.
Of 670 possible sessions from
October 2021 through May 2024, the fixed-sample probes evaluate 277.
Of these, 42 have two active legs and 123 have neither. The mean open-to-close
return is $-10.15$ basis points per evaluated session before trading
costs, with a 21-session block interval of [$-25.2$, $+4.2$] basis
points conditional on the saved predictions.

The recovered Business Standard first-publication format adds 667
distinct firm-sessions and raises the number of evaluated sessions from 277 to
340, including 110 with both legs. The new observations enter training
as well as evaluation, and the refitted mean remains negative at
$-7.73$ basis points. Table~8 of \citet{lopez2025v6} reports $+11$
basis points of mean daily drift for MiniLM headlines in a broader
RavenPack and CRSP sample. Recovering the clocks improves the ability
to evaluate both portfolio legs on the archived sample. The difference
from the reference estimate also reflects changes to texts, the
issuer universe, representation and price coverage. On the fixed
sample, the paired title and body comparisons leave substantial
uncertainty about which representation produces a higher mean drift,
and charging for entry and exit reduces the portfolio means further.
Appendix Table~\ref{tab:app-drift} reports both sensitivities.
